\subsection{Future Directions}

There are several interesting directions for extending the practicality of our work. One reason for our choice of balls-in-bins sampling is that it is more amenable to Monte Carlo accounting than Poisson sampling. Poisson sampling provides much stronger amplification than shuffling in most settings \citep{chua2024private}, so ignoring the practicality of Poisson sampling, enabling Monte Carlo accounting for correlated noise with Poisson sampling could thus enable even higher utility for private training. However, the fact that the corresponding mixture of Gaussians has $2^n$ modes is a challenging technical barrier to such a result. 

To be able to use Monte Carlo sampling to analyze the privacy amplification of the balls-in-bins sampling scheme, we need to first provide a PLD-dominating pair of distributions. In Lemma~\ref{lem:bib-dim-red}, we obtain this pair by taking the modes of $P$ to be sums of certain columns of $|\bfC|$. This is potentially too loose, and we may hope to instead use sums of columns of $\bfC$, thereby reducing the norm of each mode. In the unamplified setting, this loosening of the assumption is possible whenever $\bfC[:, i] \cdot \bfC[:, j]$ for any pair of rounds $i,j$ a user can participate in. However, amplification through sampling complicates this, and standard techniques for reducing dimensionality and adaptivity don't work anymore, see Appendix~\ref{sec:dim-red} for an example.

The number of samples we need to verify a DP guarantee using Monte Carlo accounting and e.g. Bernstein's inequality scales as $1/\delta$, which may be prohibitive for small $\delta$. It is an interesting but challenging question to determine if a more careful sampler and concentration analysis can give substantially improved sample complexity for analyzing correlated noise mechanisms. 

While we were able to implement a practical variant of balls-in-bins batching that was compatible with XLA compilation, we only used our variant to train a small-scale model. We leave it to future work to implement balls-in-bins batching at larger scales, and in particular compare the scalability of balls-in-bins to that of unamplified training and Poisson sampling.